% Tableau réponses (Série 2, gros caractères) en A3 paysage.
% Document séparé, compilé en PDF puis inséré (\includepdf) dans la feuille
% élève principale : le package geometry interdit de changer le format
% physique de la page en cours de document via \newgeometry, donc on ne
% peut pas simplement passer ces deux pages en A3 depuis le fichier
% principal (resté en A4 paysage). Le contenu ci-dessous est celui du
% tableau réponses de "Feuille élève - Série 2 ... - Gros caractères.tex",
% avec les largeurs de colonnes et les cases réponses agrandies (~x1.45)
% pour occuper la page A3 au lieu de flotter dans du vide.
\documentclass[20pt]{extarticle}
\usepackage{lmodern}
\makeatletter
\renewcommand\normalsize{\@setfontsize\normalsize{22pt}{27pt}}
\renewcommand\small{\@setfontsize\small{19pt}{24pt}}
\renewcommand\footnotesize{\@setfontsize\footnotesize{15pt}{19pt}}
\renewcommand\scriptsize{\@setfontsize\scriptsize{13pt}{15pt}}
\renewcommand\tiny{\@setfontsize\tiny{11pt}{12pt}}
\renewcommand\large{\@setfontsize\large{27pt}{33pt}}
\renewcommand\Large{\@setfontsize\Large{33pt}{39pt}}
\renewcommand\LARGE{\@setfontsize\LARGE{39pt}{45pt}}
\renewcommand\huge{\@setfontsize\huge{47pt}{57pt}}
\renewcommand\Huge{\@setfontsize\Huge{57pt}{69pt}}
\makeatother
\normalsize
\usepackage[a3paper,landscape,top=10mm,bottom=10mm,left=12mm,right=12mm]{geometry}
\usepackage[french]{babel}
\usepackage[table]{xcolor}
\usepackage{array}
\usepackage{multirow}
\usepackage{hhline}
\usepackage{ragged2e}
\usepackage{tikz}

\pagestyle{empty}
\setlength{\parindent}{0pt}
\renewcommand{\familydefault}{\sfdefault}

\definecolor{inkC}{HTML}{20232B}
\definecolor{qcmC}{HTML}{2455B8}
\definecolor{questionC}{HTML}{12805C}
\definecolor{finalC}{HTML}{B5511F}
\definecolor{emptyC}{HTML}{8A909C}

\newcommand{\qcase}[1]{%
  \begin{tikzpicture}[baseline=0pt]
    \useasboundingbox (0,0) rectangle (5.3cm,4.2cm);
    \node[anchor=north west, circle, draw=inkC, line width=1.4pt,
          inner sep=2pt, minimum size=1.3em, font=\Large, text=inkC]
      at (0.15cm,4.05cm) {#1};
  \end{tikzpicture}%
}
\newcommand{\qempty}{\cellcolor{emptyC}}

\newcolumntype{J}{>{\RaggedRight\arraybackslash}m{4.8cm}}
\newcolumntype{N}{>{\centering\arraybackslash}m{5.4cm}}
\newcolumntype{O}{>{\centering\arraybackslash}m{5.6cm}}

\begin{document}

% ======================= TABLEAU RÉPONSES (1/2 -- Jour 1-2) ==================
\noindent
Nom~: \hrulefill{} \hspace{20pt} Prénom~: \hrulefill{} \hspace{20pt} Classe~: \hrulefill

\vspace{8pt}
\renewcommand{\arraystretch}{1.3}
\setlength{\tabcolsep}{3pt}
\begin{center}
\arrayrulecolor{black}
\setlength{\arrayrulewidth}{1.4pt}
\begin{tabular}{|J|N|O|O|O|O|O|}
\hhline{-------}
 & &
\shortstack{{\LARGE\textcolor{qcmC}{A3}}\\[4pt]{\normalsize Puissances}} &
\shortstack{{\LARGE\textcolor{qcmC}{A4}}\\[4pt]{\footnotesize Écritures d'un nombre}} &
\shortstack{{\LARGE\textcolor{questionC}{Vecteurs}}\\[4pt]\parbox{5.3cm}{\centering\normalsize Égalité, somme, translations}} &
\shortstack{\parbox{5cm}{\centering\Large\textcolor{qcmC}{A9}}\\[4pt]{\normalsize Développer}\\[2pt]{\normalsize Factoriser}} &
{\LARGE\textcolor{finalC}{Algo}} \\
\hhline{-------}

\multirow{2}{4.8cm}{\textbf{Jour 1}\\[20pt] Date~:\\[15pt] \hrulefill}
  & \textcolor{qcmC}{\textbf{QCM}}
  & \qempty & \qempty & \qcase{1} & \qcase{2} & \qempty \\
\hhline{~------}
  & \textcolor{questionC}{\textbf{Questions}}
  & \qcase{3} & \qempty & \qcase{4} & \qcase{5} & \qempty \\
\hhline{-------}

\multirow{2}{4.8cm}{\textbf{Jour 2}\\[20pt] Date~:\\[15pt] \hrulefill}
  & \textcolor{qcmC}{\textbf{QCM}}
  & \qempty & \qempty & \qcase{1} & \qempty & \qcase{2} \\
\hhline{~------}
  & \textcolor{questionC}{\textbf{Questions}}
  & \qcase{3} & \qcase{4} & \qcase{5} & \qempty & \qempty \\
\hhline{-------}
\end{tabular}
\end{center}

\newpage

% ======================= TABLEAU RÉPONSES (2/2 -- Jour 3-4) ==================
\noindent
Nom~: \hrulefill{} \hspace{20pt} Prénom~: \hrulefill{} \hspace{20pt} Classe~: \hrulefill

\vspace{8pt}
\renewcommand{\arraystretch}{1.3}
\setlength{\tabcolsep}{3pt}
\begin{center}
\arrayrulecolor{black}
\setlength{\arrayrulewidth}{1.4pt}
\begin{tabular}{|J|N|O|O|O|O|O|}
\hhline{-------}
 & &
\shortstack{{\LARGE\textcolor{qcmC}{A3}}\\[4pt]{\normalsize Puissances}} &
\shortstack{{\LARGE\textcolor{qcmC}{A4}}\\[4pt]{\footnotesize Écritures d'un nombre}} &
\shortstack{{\LARGE\textcolor{questionC}{Vecteurs}}\\[4pt]\parbox{5.3cm}{\centering\normalsize Égalité, somme, translations}} &
\shortstack{\parbox{5cm}{\centering\Large\textcolor{qcmC}{A9}}\\[4pt]{\normalsize Développer}\\[2pt]{\normalsize Factoriser}} &
{\LARGE\textcolor{finalC}{Algo}} \\
\hhline{-------}

\multirow{2}{4.8cm}{\textbf{Jour 3}\\[20pt] Date~:\\[15pt] \hrulefill}
  & \textcolor{qcmC}{\textbf{QCM}}
  & \qempty & \qempty & \qcase{1} & \qcase{2} & \qempty \\
\hhline{~------}
  & \textcolor{questionC}{\textbf{Questions}}
  & \qempty & \qcase{3} & \qempty & \qempty & \qcase{4} \\
\hhline{-------}

\multirow{2}{4.8cm}{\textbf{Jour 4}\\[20pt] Date~:\\[15pt] \hrulefill}
  & \textcolor{qcmC}{\textbf{QCM}}
  & \qcase{2} & \qempty & \qcase{1} & \qempty & \qempty \\
\hhline{~------}
  & \textcolor{questionC}{\textbf{Questions}}
  & \qempty & \qcase{3} & \qempty & \qcase{4} & \qcase{5} \\
\hhline{-------}
\end{tabular}
\end{center}

\end{document}
